\section{Docker e Ambiente de Integração}

Os três repositórios permanecem independentes, mas podem ser executados em conjunto. A unidade de execução não precisa coincidir com a unidade de versionamento.

\subsection{Containers}

O ambiente integrado deverá possuir, conceitualmente, pelo menos:

\begin{itemize}
    \item um container para o frontend;
    \item um container para o Java/Spring Boot;
    \item um container para o pipeline Python;
    \item um container para o PostgreSQL.
\end{itemize}

Cada componente de aplicação pode produzir sua própria imagem Docker. Não é necessário colocar Java, Python e frontend em uma única imagem.

\placeholderfigure{docker-rede}{7cm}{Containers do ambiente integrado e comunicação pela rede Docker.}

\subsection{Comunicação entre containers}

Serviços executados na mesma rede Docker Compose podem se comunicar usando o nome do serviço como hostname. Assim, o backend Java pode acessar o PostgreSQL por um endereço conceitual como:

\begin{lstlisting}[language=YAML]
jdbc:postgresql://db:5432/sala_saude
\end{lstlisting}

O pipeline Python pode utilizar a mesma infraestrutura:

\begin{lstlisting}[language=YAML]
DB_HOST=db
DB_PORT=5432
DB_NAME=sala_saude
\end{lstlisting}

O endereço \texttt{localhost} não deve ser usado para alcançar outro container, pois dentro de um container \texttt{localhost} representa o próprio container.

\subsection{Compose de integração}

Recomenda-se que o repositório de deployment possua um Compose capaz de integrar as imagens dos três repositórios:

\begin{lstlisting}[language=YAML]
services:
  frontend:
    image: nss-frontend:dev

  java:
    image: nss-java:dev
    depends_on:
      - db

  pipeline:
    image: nss-pipeline:dev
    depends_on:
      - db

  db:
    image: postgres:16-alpine
\end{lstlisting}

Durante desenvolvimento, o Compose também pode usar \texttt{build.context} para construir as imagens a partir dos diretórios locais dos três repositórios.

\subsection{Ambientes}

A mesma separação pode ser utilizada em diferentes ambientes:

\begin{itemize}
    \item \textbf{desenvolvimento}: imagens construídas localmente;
    \item \textbf{integração}: versões coordenadas dos três componentes;
    \item \textbf{produção}: imagens versionadas publicadas em um registry.
\end{itemize}

\begin{regra}[Isolamento das imagens]
Cada componente deve possuir sua própria imagem Docker. A integração ocorre no nível de rede e orquestração, não pela fusão dos componentes em uma imagem única.
\end{regra}
